\documentclass[11pt]{article}
\usepackage[margin=2.3cm]{geometry}
\usepackage{amsmath,amssymb}
\usepackage{xr}
\externaldocument{main}
\usepackage[colorlinks=true,linkcolor=blue,urlcolor=blue]{hyperref}
\usepackage{enumitem}
\usepackage{microtype}
\setlength{\parskip}{0.45em}
\setlength{\parindent}{0pt}
\newenvironment{comment}[1]{\par\medskip\noindent\textbf{Comment~#1.}\ \itshape}{\par}
\newcommand{\resp}{\par\noindent\textbf{Response.}\ }

\title{Response to the reviewer (sixth round)\\[0.3em]\large A reliability-constrained, multi-objective
framework for minimum-mass design of steel frame structures under material and loading uncertainty}
\author{[Author name]}
\date{}

\begin{document}
\maketitle

We thank the reviewer for pointing out that our fifth-round letter did not respond to the comment
on marginal versus simultaneous confidence statements. We apologise for the omission; the comment is
answered here. The responses to Comments~2--4 in the fifth-round letter are unchanged. Apart from the
new simultaneous bounds below, which are computed from the existing verification samples, no
numerical results have changed. Section and table numbers refer to the revised manuscript.

\begin{comment}{1 (marginal versus simultaneous confidence statements)}
The one-sided 95\% lower bounds are marginal bounds for individual limit states. Insert the suggested
``Scope of confidence statements'' paragraph after the description of RBDO-C95, and make the
Abstract and Conclusions consistent. The statement that all seven RBDO-C95 designs meet both
targets with 95\% confidence should be qualified as a set of separate marginal assessments or
supported by a simultaneous confidence procedure.
\end{comment}
\resp We did both.
\begin{itemize}[itemsep=0.15em]
\item \emph{Scope paragraph.} The suggested paragraph is inserted at the end of Section~\ref{sec:mc},
immediately after the description of RBDO-C95, with the reviewer's wording. We added one sentence:
where the paper makes a simultaneous statement, it uses the Bonferroni procedure. Each of $m$ bounds
is computed at the one-sided level $1-0.05/m$, so that all $m$ hold together with probability of at
least 95\% whatever the dependence between the estimates. The definition of claim~(L) now also says
that it is a marginal statement about one limit state of one design.
\item \emph{Simultaneous procedure (tool).} We added \texttt{beta\_lcb\_simultaneous} to
\texttt{rbo/reliability.py} and the table generator \texttt{tab\_simultaneous} to
\texttt{make\_figures.py}. The latter writes \texttt{results/simultaneous.json} and the new
Table~\ref{tab:simultaneous}. The bounds use the existing $10^6$ verification samples, so no new
simulation was needed. The results for the RBDO-C95 designs are:
\begin{itemize}[itemsep=0.1em]
\item for each design on its own ($m=3$ limit states), all bounds exceed their targets, so each design
meets its three targets simultaneously with at least 95\% confidence;
\item for all seven designs together ($m=21$), all bounds exceed their targets except the ULS bound of
the low-wind frame, which is 2.995 against a target of 3.0 (marginal bound 3.005). A simultaneous
95\% statement for all seven frames together therefore cannot be made. We report this as it is.
\end{itemize}
\item \emph{Abstract.} It now says that the point-estimate and confidence-bounded designs met
\emph{each} target with \emph{marginal} 95\% confidence in the baseline surrogate model. It adds that
each RBDO-C95 design also met its three targets simultaneously (Bonferroni), but not all seven designs
jointly.
\item \emph{Results, Discussion and Conclusions.} Every statement that a design meets ``both targets
with 95\% confidence'' now reads ``each target with marginal 95\% confidence''. This covers
Sections~\ref{sec:basedesign}, \ref{sec:targets} and~\ref{sec:strict}, the caption of
Table~\ref{tab:driftlimits}, the Discussion and Conclusions~3--4. Section~\ref{sec:strict} has a new
paragraph ``Simultaneous confidence'' with Table~\ref{tab:simultaneous}. Conclusion~4 states the
per-design simultaneous result and the exception for the seven-design family. The statement that the
bounds address Monte Carlo sampling uncertainty only, and not model-form or systematic approximation
errors, is retained in each place.
\end{itemize}
The marginal and simultaneous bounds differ by less than 0.03 in $\beta_{ULS}$ and less than 0.01 in
$\beta_{drift}$. The distinction therefore changes a conclusion only for designs within a few
hundredths of a target, as the low-wind frame shows.

\section*{Remaining item}
The author and affiliation placeholders remain, because the names and affiliations have not yet
been supplied to us. They will be inserted before submission. The manuscript keeps the exact
repository commit that produced it.

\end{document}
